% Clase 2 — Ética en IA y Aplicaciones Futuras
\section{Ética en IA y Aplicaciones Futuras}

\subsection{Riesgos y Desafíos Éticos}

\begin{itemize}
    \item \textbf{Sesgo algorítmico}: sistemas entrenados con datos sesgados perpetúan o amplifican discriminación.
    \item \textbf{Privacidad}: recolección masiva de datos personales.
    \item \textbf{Seguridad}: ataques adversariales, manipulación de sistemas autónomos.
    \item \textbf{Desplazamiento laboral}: automatización de tareas cognitivas y manuales.
    \item \textbf{Concentración de poder}: dominio de pocas empresas sobre infraestructura de IA.
\end{itemize}

\subsection{Marcos Regulatorios}

\begin{itemize}
    \item \textbf{AI Act} (Unión Europea): regulación por niveles de riesgo.
    \item Lineamientos de uso de IA generativa en contextos académicos.
\end{itemize}

\subsection{Niveles de Uso de IA Generativa (Uniandes)}

\begin{tabular}{cll}
\toprule
\textbf{Nivel} & \textbf{Uso} & \textbf{Descripción} \\
\midrule
1 & Sin IA generativa & No debe usarse en ningún aspecto \\
2 & Generar/estructurar ideas & No incluir material generado en el trabajo \\
3 & Editar y refinar & Explicar uso en nota al pie o anexo \\
4 & Completar tareas (revisión humana) & Distinguir material IA; incluir transcripción de prompts \\
5 & Uso pleno & Sin restricciones de distinción \\
\bottomrule
\end{tabular}

\subsection{Futuro de la IA}

% Agrega aquí reflexiones y tendencias discutidas en clase
